%% IGVC 2026 AutoNav -- Autonomy + Cyber Security (preliminary oral). Aligned to the design report.
%% Compile: pdflatex IGVC2026_Autonomy.tex (run twice).
\documentclass[aspectratio=169,11pt]{beamer}
\usepackage[T1]{fontenc}\usepackage[utf8]{inputenc}
\usepackage[scaled=0.92]{helvet}\renewcommand{\familydefault}{\sfdefault}
\usepackage{textcomp}\usepackage{booktabs}\usepackage{graphicx}\usepackage{array}
\graphicspath{{assets/}}

\useinnertheme{default}\useoutertheme{default}\setbeamertemplate{navigation symbols}{}
\definecolor{accent}{HTML}{1F3A5F}\definecolor{accent2}{HTML}{2E6B8F}
\definecolor{muted}{HTML}{5A6672}\definecolor{ink}{HTML}{1A1A1A}
\setbeamercolor{normal text}{fg=ink}\setbeamercolor{frametitle}{fg=accent}
\setbeamercolor{itemize item}{fg=accent}\setbeamercolor{structure}{fg=accent}
\setbeamerfont{frametitle}{series=\bfseries,size=\large}
\setbeamertemplate{itemize item}{\color{accent}\rule[0.12ex]{0.6ex}{0.6ex}}
\setbeamertemplate{frametitle}{%
  \vspace{0.6em}{\usebeamerfont{frametitle}\usebeamercolor[fg]{frametitle}\insertframetitle}\par
  \vspace{0.22em}{\color{accent}\hrule height 0.9pt}\vspace{0.18em}}
\setbeamertemplate{footline}{%
  \hbox to \paperwidth{\hskip1.2em\color{muted}\scriptsize
    Cal Poly Pomona \textbullet\ IGVC 2026 AutoNav \textbullet\ MARVIN \textbullet\ Autonomy\hfil
    \color{muted}\scriptsize\insertframenumber/\inserttotalframenumber\hskip1.2em}\vskip0.5em}
\setbeamersize{text margin left=1.3em,text margin right=1.3em}

\newcommand{\imgfb}[4]{%
  \IfFileExists{assets/#1}{\includegraphics[#3]{#1}}{%
    \IfFileExists{assets/#2}{\includegraphics[#3]{#2}}{%
      \setlength{\fboxsep}{6pt}\fcolorbox{muted}{black!2}{%
        \parbox[c][3.0cm][c]{0.9\linewidth}{\centering\scriptsize\color{muted}#4}}}}}
\newcommand{\req}[1]{\vspace{0.35em}{\footnotesize\color{accent2}\textbf{Requirement} #1}}

\begin{document}

% ===== 1. PERCEPTION =====
\begin{frame}{Perception: sensing lanes and obstacles}
\begin{columns}[T]
\begin{column}{0.50\textwidth}
\begin{itemize}\setlength\itemsep{3pt}
  \item \textbf{LiDAR} (Velodyne VLP-16, 10\,Hz) --- physical obstacles in 3-D; any color, any lighting
  \item \textbf{Vision} (ZED~X) --- three color classes: white lanes, barrels, and white potholes
  \item White potholes \emph{separated} from white lane lines --- one avoided, one followed
  \item Startup auto-calibration adapts to the day's lighting; fused into one Nav2 costmap
\end{itemize}
\req{(IGVC \S II.2): detect obstacles $\geq$2\,m $\rightarrow$ \textbf{15\,m}; vision $\geq$10\,Hz $\rightarrow$ \textbf{10\,Hz}.}
\end{column}
\begin{column}{0.48\textwidth}\centering
  \imgfb{segmentation.png}{zed_rgb.jpg}{width=\linewidth,height=0.44\textheight,keepaspectratio}{Live segmentation + point-cloud costmap}\\[0.15em]
  {\tiny\color{muted}live lane / barrel / pothole segmentation $\rightarrow$ point-cloud costmap}\\[0.35em]
  \begin{minipage}{0.46\linewidth}\centering
    \imgfb{velodyne_vlp16.jpg}{x}{height=1.5cm,keepaspectratio}{VLP-16}\\{\tiny\color{muted}VLP-16 LiDAR}
  \end{minipage}\hfill
  \begin{minipage}{0.46\linewidth}\centering
    \imgfb{zed_camera.jpg}{x}{height=1.5cm,keepaspectratio}{ZED X}\\{\tiny\color{muted}ZED~X camera}
  \end{minipage}
\end{column}
\end{columns}
\note{Two sensors fuse into one Nav2 costmap. The Velodyne VLP-16 LiDAR gives 3-D obstacle geometry at ten hertz, independent of color and lighting. The ZED X camera segments three classes: lane lines, barrels, and solid-white potholes. Potholes and lane lines are both white, so we separate them at the pixel level -- potholes avoided, lanes tracked. Startup calibration adapts the vision thresholds to ambient lighting. Required obstacle detection: two meters; measured, fifteen, at full camera rate.}
\end{frame}

% ===== 2. DRIVING LOGIC =====
\begin{frame}{Driving logic: plan, then re-plan continuously}
\begin{columns}[T]
\begin{column}{0.52\textwidth}
\begin{itemize}\setlength\itemsep{3pt}
  \item \textbf{Planner} --- Navfn (Dijkstra): chosen over Smac-Hybrid, whose 2.31\,m turning circle can't fit a 10-ft lane; our tracked chassis turns in place
  \item \textbf{Controller} --- MPPI \emph{generates} trajectories at 20\,Hz to deviate around new barrels, not just follow the plan
  \item Handles potholes, switchbacks, dead-ends and the ramp; recovery tree under a 45\,s watchdog (60\,s hold-up rule)
  \item Dual EKF pose; \texttt{mission\_manager} fires GPS waypoints in sequence
\end{itemize}
\req{obstacle reaction $\leq$250\,ms $\rightarrow$ \textbf{47\,ms}; waypoint accuracy $\leq$1\,m; avg speed $\geq$1\,mph $\rightarrow$ \textbf{1.7\,mph}.}
\end{column}
\begin{column}{0.46\textwidth}\centering
  \imgfb{nav2_costmap.png}{sim_nav.png}{width=\linewidth,height=0.58\textheight,keepaspectratio}{Nav2 costmap + planned path}\\[0.25em]
  {\tiny\color{muted}Nav2 local costmap: lethal cells (cyan) ringed by the inflation cost-gradient}
\end{column}
\end{columns}
\note{Planning and control are separate. Global planner: Navfn, a Dijkstra planner -- over Smac-Hybrid, whose two-point-three-meter turning radius won't fit a ten-foot lane; the tracked chassis turns in place, so curvature constraints add nothing. Controller: MPPI, a sampling-based MPC at twenty hertz -- it re-plans around newly placed barrels and rejects any rollout that enters a pothole. A dual EKF supplies pose; a recovery behavior tree under a forty-five-second watchdog handles dead-ends and switchbacks. Required obstacle reaction: two hundred fifty milliseconds; measured, forty-seven.}
\end{frame}

% ===== 3. VALIDATION + KPIs =====
\begin{frame}{Validation: key performance indicators}
\small Two-tier testing --- a \textbf{Webots} simulation validates the full Nav2 pipeline; the \textbf{hardware} chassis proves the numbers. KPIs are set against the IGVC rules.
\vspace{0.4em}
\begin{columns}[c]
\begin{column}{0.37\textwidth}\centering
  \imgfb{sim_nav.png}{x}{width=\linewidth,height=0.4\textheight,keepaspectratio}{Webots simulation}\\[0.15em]
  {\tiny\color{muted}Webots: full-pipeline campus simulation}
\end{column}
\begin{column}{0.61\textwidth}
\renewcommand{\arraystretch}{1.25}\footnotesize
\begin{tabular}{@{}lcc@{}}
\toprule
\textbf{KPI (rule source)} & \textbf{Target} & \textbf{Measured}\\
\midrule
Obstacle reaction time & $\leq$250\,ms & \textbf{47\,ms}\\
Min.\ average speed (hold-up rule) & $\geq$1\,mph & \textbf{1.7\,mph}\\
Waypoint arrival accuracy & $\leq$1\,m & \textbf{$<$1\,m}\\
Obstacle detection range & $\geq$2\,m & \textbf{15\,m}\\
Full-course completion & $\leq$5\,min & field run pending\\
\bottomrule
\end{tabular}
\end{column}
\end{columns}
\vspace{0.45em}
\centering{\color{accent}\textbf{Field test 2026-05-21: 7 recoveries, 0 collisions around a person in the path.}}
\note{Two-tier validation: a Webots simulation runs the full Nav2 pipeline; the hardware chassis confirms the numbers. Each KPI maps to an IGVC rule, target versus measured. Obstacle reaction: forty-seven milliseconds against a two-hundred-fifty target. May twenty-first field test, person stepping into the path: seven recovery behaviors, zero collisions.}
\end{frame}

% ===== 4. LOCALIZATION =====
\begin{frame}{Localization: control in the drift-free frame}
\begin{columns}[T]
\begin{column}{0.54\textwidth}
\begin{itemize}\setlength\itemsep{3pt}
  \item Dual EKF at 30\,Hz: local frame fuses IMU + wheel odometry + ZED yaw; global frame adds GPS
  \item GPS is \textbf{unaided} (2--5\,m), fused \textbf{differentially} behind a 6$\sigma$ outlier gate --- \textbf{no RTK} (IGVC \S I.2)
  \item Per-step GPS deltas, not absolute fixes, stop $\pm$2--5\,m noise from jolting the costmap
  \item Safety-critical control runs in the \textbf{drift-free odometry frame}
\end{itemize}
\vspace{0.5em}
{\footnotesize\color{accent}\textbf{A deliberate design choice, not a workaround.}}
\end{column}
\begin{column}{0.42\textwidth}\centering
  \imgfb{gps_nav.png}{zed_rviz.jpg}{width=\linewidth,height=0.55\textheight,keepaspectratio}{GPS / odometry navigation}\\[0.25em]
  {\tiny\color{muted}GPS waypoints in the map frame; control in the drift-free odom frame}
\end{column}
\end{columns}
\note{A dual EKF at thirty hertz: the local filter fuses IMU, wheel odometry, and ZED yaw; the global filter adds GPS. No RTK -- IGVC rule I-point-two forbids positioning base stations -- so GPS is unaided two-to-five-meter SBAS, fused differentially behind a six-sigma outlier gate. Per-step GPS deltas, not absolute fixes, keep position noise out of the costmap. Safety-critical control runs in the drift-free odometry frame; GPS only biases the long-range waypoints.}
\end{frame}

% ===== 5. CYBER SECURITY =====
\begin{frame}{Cyber security: three risks, defense in depth}
\small We ran a NIST risk assessment; per IGVC \S11 here are our top three vulnerabilities.
\vspace{0.4em}
\renewcommand{\arraystretch}{1.2}\footnotesize
\begin{tabular}{@{}p{0.20\textwidth} p{0.40\textwidth} p{0.34\textwidth}@{}}
\toprule
\textbf{Vulnerability} & \textbf{What we do today} & \textbf{Production hardening}\\
\midrule
Onboard network (default DDS) & Wi-Fi limited to team laptops; judges' wireless E-stop cuts motor power & ROS\,2 DDS security (per-node certs, encryption) + VLAN / WPA3\\
Web-UI hijack (no login) & 500\,ms watchdog zeros drive on lost input; manual override; E-stop & TLS client certs + per-session tokens; Foxglove bound to localhost\\
GPS spoofing & 6$\sigma$ EKF rejects implausible jumps; GPS used as per-step correction only & RAIM multi-band receiver; authenticated RTK; dead-reckoning cross-check\\
\bottomrule
\end{tabular}
\vspace{0.5em}
\centering{\color{accent}\textbf{Defense in depth:} two independent \emph{hardware} stops --- a physical E-stop button and the judges' wireless relay --- plus a software interlock; no single software compromise can move the vehicle.}
\note{A NIST risk assessment, three top vulnerabilities. First, the network: ROS 2's default DDS has no authentication, so any host on the Wi-Fi could publish a drive command. We restrict the network to team hardware; the judges' wireless E-stop cuts motor power independently of software. Second, the web operator interface has no login -- a five-hundred-millisecond watchdog zeros the drive when commands stop, with manual override. Third, GPS spoofing: the filter takes GPS only as a bounded per-step correction and rejects physically impossible jumps. Layered defense: two hardware stops -- the physical E-stop and the judges' wireless relay -- plus a software interlock. No single software compromise can move the vehicle.}
\end{frame}

\end{document}
